\documentclass[11pt]{article}
\usepackage{amsmath,amssymb,amsthm,mathtools}
\usepackage[margin=1in]{geometry}
\usepackage{xcolor}
\usepackage[hidelinks]{hyperref}
\hypersetup{pdftitle={The eigenvalue decay of the limiting kernel of a sparse Mertens matrix},pdfauthor={Jeffery Kline},pdfsubject={Release candidate, version 0.1.0}}
\usepackage[expansion=false]{microtype}
\allowdisplaybreaks
\newtheorem{theorem}{Theorem}[section]
\newtheorem{proposition}[theorem]{Proposition}
\newtheorem{lemma}[theorem]{Lemma}
\newtheorem{corollary}[theorem]{Corollary}
\theoremstyle{definition}
\newtheorem{definition}[theorem]{Definition}
\newtheorem{example}[theorem]{Example}
\newtheorem{question}[theorem]{Question}
\theoremstyle{remark}
\newtheorem{remark}[theorem]{Remark}
\newcommand{\N}{\mathbb N}
\newcommand{\one}{\mathbf 1}
\newcommand{\sfq}{\mathcal S}
\newcommand{\cyl}{\mathcal T}
\newcommand{\rankone}[1]{R_{#1}}
\newcommand{\ncount}{\mathsf n}
\newcommand{\Tr}{\operatorname{Tr}}
\newcommand{\diag}{\operatorname{diag}}
\newcommand{\supp}{\operatorname{supp}}
\newcommand{\HS}{\mathrm{HS}}
\newcommand{\norm}[1]{\left\lVert #1\right\rVert}
\newcommand{\abs}[1]{\left\lvert #1\right\rvert}
\newcommand{\inner}[2]{\left\langle #1,#2\right\rangle}
\newcommand{\muvec}{\boldsymbol{\mu}}
\title{The eigenvalue decay of the limiting kernel\\
of a sparse Mertens matrix}
\author{Jeffery Kline}
\date{October 23, 2026\\[3pt]\small Version 0.1.0 (release candidate)}
\begin{document}
\maketitle
\begin{abstract}
Order the squarefree integers as a rooted tree in which the parent of
$d>1$ is $d/P^+(d)$, where $P^+(d)$ is the largest prime factor of $d$.
On this tree consider the symmetric kernel $K$ whose entry at an ancestor
$j$ and a descendant $k$ is $\mu(k/j)\delta_k$, where
$\delta_k=(6/\pi^2)k^{-1}\prod_{p\le P^+(k)}p/(p+1)$. It arises as the
limit of normalized Gram matrices of inverses of sparse $(0,1)$ matrices
whose bordered versions have determinant equal to the Mertens function.
An earlier release by the author identified each fixed eigenvalue of $K$
as a limit of finite singular values; here we determine how the
eigenvalues decay. We prove that $K$ is positive and compact with
\[
 \lambda_r(K)\sim\frac{6/\pi^2}{r\log r}\qquad(r\to\infty),
\]
equivalently that the number of eigenvalues exceeding $s$ is asymptotic
to $(6/\pi^2)/(s\log(1/s))$. The same law holds for a rank-one deflation
of $K$ whose eigenvalues determine, index by index, the limits of the
rescaled lower singular values of the bordered matrices, apart from the
smallest one.
The proof removes the M\"obius signs by a diagonal unitary, realizes the
kernel as a Gram matrix of nested cylinder sets for a random set of
primes, and shows that the eigenvalue count is asymptotic to the number
of cylinders of mass above the threshold. That number is evaluated with
Mertens' theorem and the Dickman function; the constant comes from the
identity $\int_0^\infty\rho=e^\gamma$. No convergence rate is obtained,
and no estimate for the Mertens function follows.
\end{abstract}
\tableofcontents
\section{Introduction}\label{sec:intro}
All logarithms are natural. Write $\sfq$ for the positive squarefree integers,
$\mu$ for the M\"obius function, and $P^+(d)$ for the largest prime factor of
$d$, with $P^+(1)=1$. Put
\[
 c_0=\frac{6}{\pi^2},\qquad
 \nu_d=\frac1d\prod_{p\le P^+(d)}\frac{p}{p+1},\qquad
 \delta_d=c_0\nu_d\quad(d\in\sfq).
\]
Thus $\nu_1=1$. The primes in each product are unrestricted by divisibility.

The squarefree integers form a rooted tree: the root is $1$, and the parent
of $d>1$ is $d/P^+(d)$. Write $j\preceq k$ when $j$ is an ancestor of $k$,
including equality. Equivalently, $j\mid k$ and every prime divisor of $k/j$
exceeds $P^+(j)$. Listing the prime factors of $k$ in increasing order, the
ancestors of $k$ are exactly the products of its initial segments.

Define a real symmetric matrix on $\N\times\N$ by
\begin{equation}\label{eq:K}
 K_{jk}=
 \begin{cases}
 \mu(k/j)\delta_k,&j,k\in\sfq,\ j\preceq k,\\
 \mu(j/k)\delta_j,&j,k\in\sfq,\ k\preceq j,\\
 0,&\text{otherwise}.
 \end{cases}
\end{equation}
The two applicable cases agree on the diagonal. Nonsquarefree coordinates
are zero. We shall show that this matrix defines a positive compact operator
on $\ell^2(\N)$.

For a positive compact operator $A$, let
\[
 \ncount(A;t)=\#\{r:\lambda_r(A)>t\},\qquad t>0,
\]
where positive eigenvalues are ordered decreasingly, with multiplicity.
All counting functions in the operator arguments use this strict convention.
The subject of this paper is the rate at which the eigenvalues of $K$ tend
to zero.

\begin{theorem}\label{thm:main}
The operator $K$ in \eqref{eq:K} is positive and compact. As $s\downarrow0$,
\begin{equation}\label{eq:maincount}
 \ncount(K;s)\sim\frac{c_0}{s\log(1/s)}.
\end{equation}
Equivalently, as $r\to\infty$,
\begin{equation}\label{eq:maineigen}
 \lambda_r(K)\sim\frac{c_0}{r\log r}.
\end{equation}
Consequently,
\[
 \sum_{r\le R}\lambda_r(K)\sim c_0\log\log R.
\]
The same three asymptotics hold for the positive compact operator
\begin{equation}\label{eq:deflation}
 L=K-\frac{(Ke_1)(Ke_1)^*}{c_0}.
\end{equation}
\end{theorem}

Two immediate consequences are that $K$ and $L$ lie in the Schatten class
$\mathcal S_p$ exactly when $p>1$ (Corollary~\ref{cor:schatten}), and that
$K$ is not trace class even though $K$ minus its diagonal is
(Proposition~\ref{prop:traceclass}).

\subsection{Where the kernel comes from}
The kernel $K$ appeared in the author's earlier study~\cite{Kline} of a
sparse matrix, introduced in~\cite{Kline19}, whose determinant is the
Mertens function $M(n)=\sum_{m\le n}\mu(m)$. Let $A_n$ be
the $n\times n$ matrix with ones on the diagonal and an additional one in
position $(i,i/P^+(i))$ for each squarefree $i$ with $1<i\le n$. Its inverse
$C_n=A_n^{-1}$ has entries $\mu(i/j)$ when $j\preceq i$
(Section~\ref{sec:finite}). Replacing the first row of $A_n$ by ones gives
a matrix $B_n$ with $\det B_n=M(n)$. The following facts are proved
in~\cite{Kline}.
\begin{itemize}
\item \cite[Theorem~4.3]{Kline}: the matrices $C_n^*C_n/n$, extended by zero,
 converge to $K$ in Hilbert--Schmidt norm. Hence, for every fixed $r$,
 $s_r(C_n)/\sqrt n\to\sqrt{\lambda_r(K)}$.
\item \cite[Theorem~4.6]{Kline}: for every fixed $r\ge1$,
 $\sqrt n\,\sigma_{n-r}(B_n)\to\lambda_r(L)^{-1/2}$, with $L$ as in
 \eqref{eq:deflation}.
\end{itemize}
Both statements fix the index $r$ and let $n\to\infty$. They identify each
limit but do not describe how the limits behave as $r$ grows. The earlier
analysis also gives little quantitative information about the limiting
operator itself: Hilbert--Schmidt membership gives
$\sum_r\lambda_r(K)^2<\infty$, but not the order of $\lambda_r(K)$.
Theorem~\ref{thm:main} answers this question for the limiting operators.
It does not, by itself, give information about $s_{r}(C_n)$ when $r$ and
$n$ grow together; Section~\ref{sec:finite} records what an operator-norm
error estimate would add.

A first guess comes from the diagonal. The diagonal entries of $K$ are
$\delta_d$, and Mertens' product theorem gives
$\delta_d\approx e^{-\gamma}/(d\log P^+(d))$ when $P^+(d)$ is large. The
off-diagonal part is not small in the sense needed, however. It is trace
class (Proposition~\ref{prop:traceclass}), which controls partial sums of
eigenvalues up to a bounded error. The eigenvalues themselves, however,
are of order $1/(r\log r)$, and a trace-class operator may have singular
values that are not $o(1/(r\log r))$. Trace-class membership alone
therefore does not transfer an asymptotic for individual eigenvalues from
the diagonal to $K$. The proof uses the structure of $K$ instead.

\subsection{Outline of the proof}
The argument has four steps.

\emph{Removing the signs.} On comparable squarefree indices
$\mu(k/j)=\mu(j)\mu(k)$, so conjugation by the diagonal unitary with
entries $\mu(d)$ replaces $K$ by the unsigned kernel with entries
$\delta_{\max(j,k)}$ on comparable pairs. The spectrum is unchanged.

\emph{A cylinder representation.} Choose a random set of primes, including
each prime $p$ independently with probability $1/(p+1)$. For $d\in\sfq$ let
$A_d$ be the event that the selected primes up to $P^+(d)$ are exactly the
prime divisors of $d$. Then $A_d$ has probability $\nu_d$; the events are
nested along the tree and disjoint for incomparable indices. The unsigned
kernel is therefore $c_0$ times the Gram matrix of the indicators
$\one_{A_d}$ (Section~\ref{sec:kernel}).

\emph{From masses to eigenvalues.} Section~\ref{sec:tree} treats any tree
of nested sets whose children are disjoint and carry at most a fixed
fraction $\theta<1$ of their parent's mass. The operator $T$, the sum
over vertices $v$ of the rank-one maps
$f\mapsto\inner{f}{\one_{A_v}}\one_{A_v}$, is compact, and splitting each subtree into
its root and its children expresses $T$ as a rank-one operator plus an
orthogonal sum of rescaled copies of the subtree operators. Rank-one
interlacing, applied recursively, traps the eigenvalue count of $T$
between the number of leaves and an inflated number of vertices of the
finite tree $\{v:\nu(A_v)>t\}$. When the child-to-parent mass ratios tend
to zero outside finite sets and the mass count is regularly varying, both
bounds are asymptotic to the mass count (Theorem~\ref{thm:transfer}).

\emph{Counting the masses.} It remains to count $d$ with $\nu_d>t$
(Section~\ref{sec:arithmetic}). Grouping $d$ by its largest prime factor
reduces this to counts of squarefree smooth numbers, which we evaluate
with the Dickman function and Mertens' theorems. The answer is
$D(t)\sim1/(t\log(1/t))$. Its constant is the product of the squarefree
density $c_0$, the limit $\zeta(2)e^{-\gamma}$ of $m(p)\log p$, where
$m(p)=\prod_{q\le p}q/(q+1)$, and the classical value
$\int_0^\infty\rho(v)\,dv=e^\gamma$; since $c_0\zeta(2)=1$, the product
is one. The factor $c_0$ in Theorem~\ref{thm:main} comes only from the
relation $\delta_d=c_0\nu_d$.

\subsection{Scope}
The result is unconditional. Its ingredients are classical: rank-one
interlacing, the min--max principle, smooth-number estimates
from~\cite{HT}, and Mertens' theorems. What this paper contributes is the
cylinder representation of $K$, the reduction of its eigenvalue count to a
count of cylinder masses, and the evaluation of that count, which together
give the exact asymptotic in Theorem~\ref{thm:main}. Section~\ref{sec:related}
discusses related results on arithmetic matrices, operators on trees, and
self-similar spectral problems, and lists comparisons that remain open.

We do not obtain a rate of convergence or a second-order term in
\eqref{eq:maincount}. The theorem concerns the limiting operators; it
does not give a growing-index statement for the finite matrices without
an additional estimate (Corollary~\ref{cor:growing}). The kernel is
positive after a diagonal change of signs, and the proof uses no
cancellation in sums of $\mu$; in particular, no new estimate for $M(n)$
follows.

Appendix~\ref{sec:profiles} gives a second sufficient condition for the
transfer from masses to eigenvalues, based on an exact rank-one shift
identity, and Appendix~\ref{sec:stationary} applies it to stationary
trees and computes the spectrum of homogeneous trees. Neither appendix is
needed for Theorem~\ref{thm:main}. Whether the transfer requires an
assumption beyond regular variation of the mass count is left open
(Question~\ref{q:drop}).

\section{The cylinder representation}\label{sec:kernel}
The proof separates the geometry of nested sets from an arithmetic counting
problem. The relevant diagonal count is
\begin{equation}\label{eq:Darith}
 D(t)=\#\{d\in\sfq:\nu_d>t\}.
\end{equation}
Section~\ref{sec:tree} gives an abstract criterion for transferring $D$ to
the spectrum. Section~\ref{sec:arithmetic} proves
$D(t)\sim[t\log(1/t)]^{-1}$, including the leading constant. In this
section we realize $K$, up to a diagonal unitary and the factor $c_0$, as a
Gram matrix of nested sets whose masses are the numbers $\nu_d$.

\subsection{A product probability space}
Let $\Omega=\{0,1\}^{\{\text{primes}\}}$, with independent coordinates and
probability $1/(p+1)$ that prime $p$ is selected. Denote its probability
measure by $\nu$. This is the local model of a random squarefree integer:
$1/(p+1)$ is the density, among squarefree integers, of those divisible
by $p$. For $d\in\sfq$, let
\[
 A_d=\{\omega:\text{the selected primes at most }P^+(d)
                   \text{ are exactly the prime divisors of }d\};
 \qquad A_1=\Omega.
\]
Direct multiplication gives
\begin{equation}\label{eq:cylinder-mass}
 \nu(A_d)=
 \prod_{p\mid d}\frac1{p+1}
 \prod_{\substack{p\le P^+(d)\\p\nmid d}}\frac p{p+1}
 =\nu_d.
\end{equation}
If $j\preceq k$, then $A_k\subseteq A_j$; incomparable indices give disjoint
cylinders. Indeed, if $P^+(j)\le P^+(k)$ and the two cylinders intersect,
the selected primes at most $P^+(j)$ force $j\preceq k$. Hence
\begin{equation}\label{eq:gramentries}
 \inner{\one_{A_j}}{\one_{A_k}}
 =\nu_{\max(j,k)}\one_{\{j,k\text{ comparable}\}}.
\end{equation}

The children of $d$ are $dq$ for primes $q>P^+(d)$. Their cylinders cover
$A_d$ up to a null set, since
$\prod_{p>P^+(d)}p/(p+1)=0$. Their relative masses are
\begin{equation}\label{eq:arith-ratios}
 \frac{\nu_{dq}}{\nu_d}
 =\frac1{q+1}\prod_{P^+(d)<p<q}\frac p{p+1}
 \le\frac1{P^+(d)+2}\le\frac13.
\end{equation}
For every fixed $y$ there are only finitely many $d\in\sfq$ with
$P^+(d)\le y$. The largest immediate-child ratio therefore tends to zero
outside finite sets of vertices, in the sense of
Definition~\ref{def:regvar} below. Also $D(t)<\infty$, since $\nu_d\le1/d$.

\subsection{Removing the signs}
Let $S$ be the diagonal unitary whose $d$th entry is $\mu(d)$ on squarefree
coordinates and $1$ otherwise. On comparable squarefree indices,
$\mu(k/j)=\mu(j)\mu(k)$. Equations \eqref{eq:K} and \eqref{eq:gramentries}
therefore identify
\begin{equation}\label{eq:gram-identification}
 K=c_0 S V^*VS,\qquad
 Vx=\sum_{d\in\sfq}x_d\one_{A_d},
 \qquad T=VV^*.
\end{equation}
Here $V$ is zero on nonsquarefree coordinates; its bounded extension and
compactness follow from Proposition~\ref{prop:compact}. Since the nonzero
spectra of $V^*V$ and $VV^*$ coincide, with multiplicity,
\begin{equation}\label{eq:countscale}
 \ncount(K;s)=\ncount(T;s/c_0).
\end{equation}
The operator $T$ acts on $L^2(\nu)$ and is the sum over $d\in\sfq$ of the
rank-one operators $f\mapsto\inner{f}{\one_{A_d}}\one_{A_d}$. The next
section studies such sums for general trees of nested sets.

\section{Compact cylinder operators and a counting theorem}\label{sec:tree}
This section is independent of arithmetic. It shows that a sum of rank-one
operators attached to a tree of nested sets is compact, bounds its
eigenvalue count above and below by counts of vertices, and gives conditions
under which both bounds are asymptotic to the number of sets of mass above
the threshold. Section~\ref{sec:consequences} applies it to the tree of
Section~\ref{sec:kernel}.

Let $\cyl$ be a countable rooted tree, each vertex having finite depth.
On a probability space $(\Omega,\nu)$, associate to each vertex $v$ a
measurable set $A_v$ of positive mass $m_v=\nu(A_v)$, with $A_o=\Omega$ at
the root. Child sets lie inside the parent and are pairwise disjoint modulo
null sets. They need not cover the parent. Set
\[
 a_v=\sup_{c\text{ child of }v}\frac{m_c}{m_v},
 \qquad a_v=0\text{ at a leaf},\qquad
 D(t)=\#\{v:m_v>t\}.
\]
Throughout this section assume
\begin{equation}\label{eq:tree-assumptions}
 0\le a_v\le\theta<1,\qquad 0<\theta<1,\qquad D(t)<\infty\quad(t>0),
 \qquad M_\theta=\frac1{1-\theta}.
\end{equation}
For $f$ in a Hilbert space, $R_f$ denotes the positive rank-one operator
with quadratic form $\inner{x}{R_fx}=|\inner{x}{f}|^2$.

\begin{proposition}[Construction and compactness]\label{prop:compact}
The sum
\[
 T=\sum_{v\in\cyl}R_{\one_{A_v}}
\]
defines a positive compact operator on $L^2(\nu)$, with $\norm T\le M_\theta$.
The synthesis map $V:\ell^2(\cyl)\to L^2(\nu)$ is compact and satisfies
$T=VV^*$; the Gram operator $V^*V$ has entries $\nu(A_v\cap A_w)$.
If $T_v$ is the sum over descendants of $v$, including $v$, then
\[
 \norm{T_v}\le M_\theta m_v.
\]
For the finite truncation $T^{(\varepsilon)}=\sum_{m_v>\varepsilon}
R_{\one_{A_v}}$, with $0<\varepsilon<1$,
\begin{equation}\label{eq:cutoff-error}
 \norm{T-T^{(\varepsilon)}}\le M_\theta\varepsilon.
\end{equation}
\end{proposition}
\begin{proof}
At depth $\ell$, the sets are disjoint and their masses are at most
$\theta^\ell$. Thus the sum $Q_\ell$ at that depth is an orthogonal
sum of rank-one operators of norms $m_v$, and $\norm{Q_\ell}\le\theta^\ell$.
Each $Q_\ell$ is compact: for every positive threshold only finitely many
of these norms exceed it. Hence $\sum_{\ell\ge0}Q_\ell$ converges in operator
norm to a positive compact operator, with norm at most $M_\theta$.
The same argument within $A_v$ gives the subtree bound.

The family $\{v:m_v>\varepsilon\}$ is finite and ancestor closed.
The omitted vertices are the disjoint subtrees rooted at minimal omitted
vertices $c$, all of mass at most $\varepsilon$. Their operators act on
orthogonal spaces $L^2(A_c)$, proving \eqref{eq:cutoff-error}.

For any $f$, the sum of $|\inner{f}{\one_{A_v}}|^2$ is $\inner{f}{Tf}$.
This defines a bounded analysis operator $V^*$, whose adjoint is the stated
synthesis map. Then $VV^*=T$. Compactness of $T$ implies compactness of $V$,
for example by polar decomposition. The Gram entries follow on coordinate
vectors. These constructions also show that finite vertex sums converge
in norm: a finite set containing all vertices of mass above $\varepsilon$
leaves a positive remainder dominated by the one in \eqref{eq:cutoff-error}.
\end{proof}

Use the normalized measure $\nu/m_v$ on $A_v$ to define the normalized
subtree operator $\widehat T_v$. Extension after multiplication by
$m_v^{-1/2}$ gives a unitary $J_v$ onto $L^2(A_v,\nu)$, and
\[
 T_v=m_v J_v\widehat T_vJ_v^*.
\]
Decomposing into child spaces, with the unused part of $A_v$ assigned a
zero operator, gives
\begin{equation}\label{eq:normalized-decomposition}
 \widehat T_v=R_{\one}+
 \bigoplus_{c\text{ child of }v}\frac{m_c}{m_v}
 U_c\widehat T_cU_c^* .
\end{equation}
Here $U_c$ is the corresponding normalized embedding and
$\norm{\one}_{L^2(A_v,\nu/m_v)}=1$. Consequently
\begin{equation}\label{eq:subtree-norm}
 1\le\norm{\widehat T_v}\le1+M_\theta a_v\le M_\theta.
\end{equation}

\begin{proposition}[Tree sandwich]\label{prop:sandwich}
Let $\mathcal L_t$ be the leaves of the finite ancestor-closed tree
$\{v:m_v>t\}$. Then
\begin{equation}\label{eq:sandwich}
 |\mathcal L_t|\le\ncount(T;t)
 \le\#\{v:m_v\norm{\widehat T_v}>t\}.
\end{equation}
In particular,
\begin{equation}\label{eq:rough-sandwich}
 D(t)-D(t/\theta)\le\ncount(T;t)\le D((1-\theta)t).
\end{equation}
\end{proposition}
\begin{proof}
Leaf vertices are incomparable, so the compression of $V^*V$ to their
coordinates is diagonal with entries $m_v>t$. The min--max principle and
equality of the nonzero spectra of $V^*V$ and $VV^*$ prove the lower bound.

Write $n_v(u)=\ncount(\widehat T_v;u)$. Positive rank-one interlacing in
\eqref{eq:normalized-decomposition} gives
\[
 n_v(u)\le
 \one_{\{u<\norm{\widehat T_v}\}}
 +\sum_{c\text{ child of }v}n_c(um_v/m_c).
\]
For $u\ge\norm{\widehat T_v}$ both sides' child counts vanish, which explains
the indicator rather than an unrestricted $1$. For a fixed $u>0$, iteration
terminates uniformly: at depth $h$ the threshold is at least $u\theta^{-h}$,
and all normalized operator norms are at most $M_\theta$. Choose $h$ with
$u\theta^{-h}\ge M_\theta$. All contributing vertices have
$m_w>um_v/M_\theta$, so only finitely many can contribute.
Unrolling proves the upper bound at the root.

A nonleaf of $\{m_v>t\}$ has a child of mass above $t$, hence has mass
above $t/\theta$. This proves the first rough bound. The second follows
from \eqref{eq:subtree-norm}.
\end{proof}

\begin{definition}\label{def:regvar}
A positive function $D$ is regularly varying at zero with index $-\alpha$
if, for every $b>0$, $D(bt)/D(t)\to b^{-\alpha}$ as $t\downarrow0$.
We say $a_v\to0$ outside finite sets if
$\{v:a_v>\eta\}$ is finite for every $\eta>0$.
\end{definition}
The second condition is stronger than convergence of $a_v$ to zero along
each infinite path of the tree. It requires that, for each $\eta>0$, all
but finitely many vertices have every child of relative mass at most
$\eta$.

\begin{theorem}[Transfer from masses to eigenvalues]\label{thm:transfer}
Under \eqref{eq:tree-assumptions}, suppose $a_v\to0$ outside finite sets
and $D$ is regularly varying at zero with index $-\alpha$, where $\alpha>0$.
Then
\[
 \ncount(T;t)\sim D(t)\qquad(t\downarrow0).
\]
The conclusion holds as well for the nonzero spectrum of the Gram operator
$V^*V$.
\end{theorem}
\begin{proof}
Fix $0<\eta<1$ and let $F_\eta=\{v:a_v>\eta\}$.
Outside $F_\eta$, a nonleaf of $\{m_v>t\}$ has mass above $t/\eta$, and
$\norm{\widehat T_v}\le1+M_\theta\eta$. Proposition~\ref{prop:sandwich} gives
\[
 D(t)-|F_\eta|-D(t/\eta)
 \le\ncount(T;t)
 \le |F_\eta|+D\bigl(t/(1+M_\theta\eta)\bigr).
\]
Divide by $D(t)$ and let $t\downarrow0$. Since $\alpha>0$, $D(t)\to\infty$.
The lower limiting ratio is at least $1-\eta^\alpha$, and the upper
limiting ratio is at most $(1+M_\theta\eta)^\alpha$. Let $\eta\downarrow0$.
\end{proof}

\section{The arithmetic mass count}\label{sec:arithmetic}
This section evaluates the count \eqref{eq:Darith}. We write each
$d>1$ as its largest prime factor times a squarefree number with smaller
prime factors, count the latter with smooth-number estimates, and sum over
the largest prime. Because the classical estimates are uniform only in
limited ranges, the sum is split into three ranges of the largest prime.

For a prime $p$, put
\[
 m(p)=\prod_{q\le p}\frac q{q+1}.
\]
Let $\Psi(x,y)$ count the positive integers $n\le x$ with $P^+(n)\le y$.
Use $\Psi(x,p^-)$ for the strict condition $P^+(n)<p$, and let
$\Psi_{\mathrm{sf}}(x;p^-)$ impose squarefreeness as well.
All three functions are zero for $x<1$.

\begin{theorem}\label{thm:arithmetic}
For the masses $\nu_d$ of Section~\ref{sec:intro},
\[
 D(t)=\#\{d\in\sfq:\nu_d>t\}
 \sim\frac1{t\log(1/t)}\qquad(t\downarrow0).
\]
\end{theorem}

\subsection{Classical inputs and the uniformity actually used}
We use Chebyshev's upper bound $\pi(y)\ll y/\log y$ and the two Mertens
asymptotics
\[
 \prod_{p\le y}(1-1/p)\sim\frac{e^{-\gamma}}{\log y},
 \qquad \sum_{p\le y}\frac1p=\log\log y+B+o(1).
\]
In particular,
\begin{equation}\label{eq:mproduct}
 m(p)\sim\frac{\zeta(2)e^{-\gamma}}{\log p},
 \qquad \frac{c}{\log p}\le m(p)\le\frac C{\log p}\quad(p\ge2)
\end{equation}
for absolute positive constants $c,C$.

The Dickman function $\rho$ is $1$ on $[0,1]$, zero on $(-\infty,0)$, and
satisfies $u\rho'(u)=-\rho(u-1)$ for $u>1$. We use:
\begin{align}
 \Psi(x,y)&\ll x e^{-u/2}
 &&(x\ge y\ge2,\ u=\log x/\log y),\label{eq:rankin}\\
 \sup_{u\in I}\left|\frac{\Psi(x,x^{1/u})}{x}-\rho(u)\right|&\longrightarrow0
 &&(x\to\infty),\label{eq:dickmancompact}\\
 \int_0^\infty\rho(u)\,du&=e^\gamma.\label{eq:dickmanintegral}
\end{align}
In \eqref{eq:dickmancompact}, $I$ is any fixed compact interval in
$(0,\infty)$.

These facts may be sourced to Hildebrand--Tenenbaum~\cite{HT}.
Printed page~414 states the fixed-$u$ Dickman limit. Equation~(1.4) there
has an extra factor $\log y$ in \eqref{eq:rankin}; the following sentence
explicitly states that this factor can be removed, citing Tenenbaum's
Theorem~III.5.1. Thus \eqref{eq:rankin} uses that strengthened bound,
not equation~(1.4) alone. For \eqref{eq:dickmanintegral}, use Lemma~2.6
on printed page~426 and the definition of $E$ in (2.3), printed page~422:
the Laplace transform is $\exp(\gamma+E(-s))$, and $E(0)=0$.

For completeness, the compact uniformity in \eqref{eq:dickmancompact}
follows from the fixed-$u$ limit. Both
$u\mapsto\Psi(x,x^{1/u})/x$ and $\rho(u)$ are nonincreasing; $\rho$ is
continuous. A finite mesh on $I$ with small oscillation of $\rho$ brackets
all intermediate values between endpoint values, where convergence is
simultaneous. For $u<1$ the smooth count is simply $\lfloor x\rfloor$.
No approximation on a growing $u$-range is used below.

\subsection{An exact weak-count identity}
Let
\[
 F(X)=\#\{d\in\sfq:\nu_d\ge1/X\}\qquad(X>1).
\]
Every $d>1$ factors uniquely as $d=pm$ with $p=P^+(d)$ and
$m$ squarefree with $P^+(m)<p$. Thus
\begin{equation}\label{eq:Fidentity}
 F(X)=1+\sum_p\Psi_{\mathrm{sf}}(x_p;p^-),
 \qquad x_p=\frac{Xm(p)}p.
\end{equation}
Primes $p\ge X$ give zero. We prove $F(X)\sim X/\log X$, then transfer to
the strict count. Write $L=\log X$ and $\alpha_p=(\log p)/L$.
Fix $0<a\le1/4<b<1$ and split the sum at $X^a$ and $X^b$.

\subsection{Small largest primes}
The term $p=2$ equals $1$ for all sufficiently large $X$, since $m=1$
is the only squarefree positive integer with all prime factors below $2$.
For $3\le p\le X^a$, \eqref{eq:mproduct} implies, uniformly in this range,
\[
 x_p\ge \frac{cX^{1-a}}{L}\ge p,\qquad \log x_p\ge L/2
\]
once $X$ is large. Hence \eqref{eq:rankin} gives
\[
 \Psi_{\mathrm{sf}}(x_p;p^-)
 \le\Psi(x_p,p)
 \ll\frac{X}{p\log p}\exp\left(-\frac{L}{4\log p}\right).
\]
For integer $k\ge1$, Chebyshev's bound gives
\[
 \sum_{e^k<p\le e^{k+1}}\frac1p
 \le e^{-k}\pi(e^{k+1})\ll\frac1k.
\]
On such a bin, $1/\log p<1/k$ and
$e^{-L/(4\log p)}\le e^{-L/(4(k+1))}\le e^{-L/(8k)}$, since $4(k+1)\le8k$
for $k\ge1$. The small-prime contribution, apart from $p=2$, is therefore
at most
\[
 CX\sum_{1\le k\le aL}\frac1{k^2}e^{-L/(8k)}.
\]
The function $f(u)=u^{-2}e^{-L/(8u)}$ is unimodal, with maximum
$256e^{-2}/L^2$. Comparing the sum with its integral and twice its maximum,
and using
\[
 \int_0^{aL}u^{-2}e^{-L/(8u)}\,du
 =\frac8L e^{-1/(8a)},
\]
we obtain
\begin{equation}\label{eq:smallrange}
 \sum_{p\le X^a}\Psi_{\mathrm{sf}}(x_p;p^-)
 \ll \frac XL e^{-1/(8a)}+\frac X{L^2}+1.
\end{equation}
The implied constant here can be taken independent of $a$ in $(0,1/4]$.

\subsection{The middle range}
We first establish, uniformly for $X^a<p<X^b$,
\begin{equation}\label{eq:squarefree-smooth}
 \Psi_{\mathrm{sf}}(x_p;p^-)
 =x_p\{c_0\rho(u_p)+o(1)\},
 \qquad u_p=\frac{\log x_p}{\log p}.
\end{equation}
The squarefree identity $\mu^2(n)=\sum_{\beta^2\mid n}\mu(\beta)$ gives
\[
 \Psi_{\mathrm{sf}}(x;p^-)
 =\sum_{\substack{\beta\le\sqrt x\\P^+(\beta)<p}}
   \mu(\beta)\Psi(x/\beta^2,p^-).
\]
Terms with $\beta>B$ have absolute sum at most
$x\sum_{\beta>B}\beta^{-2}\ll x/B$. For fixed $B$, all $\beta\le B$
satisfy $P^+(\beta)<p$ once $X$ is large. Uniformly in the middle range,
$x_p/\beta^2\to\infty$ and
\[
 \frac{\log(x_p/\beta^2)}{\log p}
\]
stays in a fixed compact subinterval of $(0,\infty)$.
Replacing the strict cutoff $p^-$ by $p$ changes a count at argument $z$
by at most $z/p$. Apply \eqref{eq:dickmancompact}, and then continuity
of $\rho$, to obtain
\[
 \frac{\Psi(x_p/\beta^2,p^-)}{x_p}
 =\beta^{-2}\rho(u_p)+o(1)
\]
uniformly for this finite set of $\beta$. First let $X\to\infty$ and then
$B\to\infty$. Since $\sum_{\beta\ge1}\mu(\beta)/\beta^2=1/\zeta(2)=c_0$,
this proves \eqref{eq:squarefree-smooth}.

Equation~\eqref{eq:mproduct} now gives, uniformly on the middle range,
\[
 u_p=\frac{1-\alpha_p}{\alpha_p}+o(1),\qquad
 c_0x_p=\frac{e^{-\gamma}X}{p\log p}\{1+o(1)\}.
\]
The total mass $\sum_{X^a<p<X^b}x_p$ is $O_{a,b}(X/L)$, so the uniform
errors remain $o_{a,b}(X/L)$ after summation. Consequently
\begin{equation}\label{eq:middlebefore}
 \frac LX\sum_{X^a<p<X^b}\Psi_{\mathrm{sf}}(x_p;p^-)
 =e^{-\gamma}\sum_{X^a<p<X^b}
 \frac1p\,\frac{\rho((1-\alpha_p)/\alpha_p)}{\alpha_p}+o(1).
\end{equation}
Mertens' reciprocal-prime asymptotic implies weak convergence of these
prime measures on $[a,b]$ to $d\alpha/\alpha$: on any fixed subinterval
$[u,v]\subset[a,b]$ their masses tend to $\log(v/u)$.
Approximating the continuous integrand in \eqref{eq:middlebefore} by
step functions yields
\begin{equation}\label{eq:middlelimit}
 \lim_{X\to\infty}\frac LX
 \sum_{X^a<p<X^b}\Psi_{\mathrm{sf}}(x_p;p^-)
 =e^{-\gamma}\int_a^b
     \rho\!\left(\frac{1-\alpha}{\alpha}\right)\frac{d\alpha}{\alpha^2}
 =e^{-\gamma}\int_{(1-b)/b}^{(1-a)/a}\rho(v)\,dv.
\end{equation}

\subsection{Large largest primes and completion}
The trivial estimate $\Psi_{\mathrm{sf}}(x_p;p^-)\le x_p$ and
\eqref{eq:mproduct} show
\begin{equation}\label{eq:largerange}
 \sum_{X^b\le p<X}\Psi_{\mathrm{sf}}(x_p;p^-)
 \le \frac{CX}{bL}\sum_{X^b\le p<X}\frac1p
 =\frac{CX}{bL}\{\log(1/b)+o(1)\}.
\end{equation}
All summands in \eqref{eq:Fidentity} are nonnegative. Thus the limiting
middle contribution is a lower bound for $\liminf LF(X)/X$, while
\eqref{eq:smallrange} and \eqref{eq:largerange} bound the excess in the
upper limit. Let $a\downarrow0$ and $b\uparrow1$, after taking the limits
in $X$. Equation~\eqref{eq:dickmanintegral} gives
\[
 \frac{LF(X)}X\longrightarrow
 e^{-\gamma}\int_0^\infty\rho(v)\,dv=1.
\]
Finally, for every fixed $\varepsilon>0$,
\[
 F\bigl(X/(1+\varepsilon)\bigr)\le D(1/X)\le F(X).
\]
Taking $X\to\infty$ and then $\varepsilon\downarrow0$ proves
Theorem~\ref{thm:arithmetic}. In particular, $D$ is regularly varying
at zero with index $-1$.

\section{Proof of the main theorem and further consequences}\label{sec:consequences}
We now combine the cylinder representation, the transfer theorem, and the
arithmetic count.

\begin{proof}[Proof of Theorem~\ref{thm:main}]
Equations \eqref{eq:arith-ratios} and \eqref{eq:cylinder-mass} verify
\eqref{eq:tree-assumptions} with $\theta=1/3$. The child ratios tend to zero
outside finite sets. Proposition~\ref{prop:compact} and
\eqref{eq:gram-identification} establish positivity and compactness.
Theorems~\ref{thm:transfer} and \ref{thm:arithmetic}, followed by
\eqref{eq:countscale}, give
\[
 \ncount(K;s)\sim D(s/c_0)
 \sim\frac{c_0}{s\log(c_0/s)}
 \sim\frac{c_0}{s\log(1/s)}.
\]

To invert the count without assumptions about multiplicities, put
$b_r=c_0/(r\log r)$ for $r\ge2$. For each fixed $\varepsilon\in(0,1)$,
\[
 \frac{\ncount(K;(1\pm\varepsilon)b_r)}r
 \longrightarrow\frac1{1\pm\varepsilon}.
\]
For all sufficiently large $r$ this traps $\lambda_r(K)$ between
$(1-\varepsilon)b_r$ and $(1+\varepsilon)b_r$.
Letting $\varepsilon\downarrow0$ proves \eqref{eq:maineigen}.
Summation, with finitely many initial terms absorbed, gives
$\sum_{r\le R}\lambda_r(K)\sim c_0\log\log R$.

Let $P=R_{\one}$ be the orthogonal projection onto constants in $L^2(\nu)$.
The root column of $V$ is $\one$, of norm one, and $Se_1=e_1$.
Thus
\[
 L=c_0 S V^*(I-P)VS\ge0.
\]
It is compact, and $K-L$ is positive of rank one. Interlacing gives
\[
 0\le\ncount(K;s)-\ncount(L;s)\le1.
\]
This proves the same counting law for $L$; inversion and summation
apply unchanged.
\end{proof}

\begin{corollary}\label{cor:schatten}
For every $p>0$, $K$ and $L$ belong to the Schatten class $\mathcal S_p$
if and only if $p>1$. In particular, they are Hilbert--Schmidt but not
trace class.
\end{corollary}
\begin{proof}
For a positive compact operator, membership in $\mathcal S_p$ means
\[
 \sum_r\lambda_r^p<\infty.
\]
Compare with $\sum_{r\ge2}(r\log r)^{-p}$, using Theorem~\ref{thm:main}.
\end{proof}

\subsection{A trace-class perturbation, and its precise use}
There is also a direct trace-class relation to the diagonal. It is useful
for sums of eigenvalues, although the pointwise counting law above uses
the tree theorem.

\begin{proposition}\label{prop:traceclass}
Let $\Delta=\diag(\delta_d:d\in\sfq)$, with zero entries elsewhere.
Then $K-\Delta$ is trace class. If $\delta_r^\downarrow$ lists the
positive diagonal entries in decreasing order, then, for every $R\ge1$,
\[
 \left|\sum_{r\le R}\lambda_r(K)
             -\sum_{r\le R}\delta_r^\downarrow\right|
 \le\norm{K-\Delta}_{\mathcal S_1}.
\]
Moreover, for each prime $p$,
\begin{equation}\label{eq:layer}
 \sum_{\substack{d\in\sfq\\P^+(d)=p}}\delta_d=\frac{c_0}{p+1}.
\end{equation}
\end{proposition}
\begin{proof}
For squarefree $j$, let $b_j$ have coordinate $K_{jk}$ at strict
descendants $k$ of $j$ and zero elsewhere. If $j>1$ and $y=P^+(j)$,
write such descendants as $k=jd$. Then $d>y$ and
$\delta_{jd}/\delta_j\le1/d$, so
\[
 \norm{b_j}^2\le\delta_j^2\sum_{d>y}d^{-2}\le\frac{\delta_j^2}{y}.
\]
The root vector $b_1$ is in $\ell^2$ since $\delta_k\le c_0/k$.
On the other hand,
\[
 \sum_{P^+(d)=p}\delta_d
 =\frac{c_0m(p)}p
       \prod_{q<p}(1+1/q)
 =\frac{c_0}{p+1},
\]
proving \eqref{eq:layer}. Hence
\[
 \sum_{\substack{j\in\sfq\\j>1}}\norm{b_j}
 \le c_0\sum_p\frac1{(p+1)\sqrt p}<\infty.
\]
Since $b_j\perp e_j$, the self-adjoint rank-two operator
$e_jb_j^*+b_je_j^*$ has trace norm $2\norm{b_j}$. Thus
\[
 K-\Delta=\sum_{j\in\sfq}(e_jb_j^*+b_je_j^*)
\]
converges in trace norm; its entries identify it with the stated
difference. The eigenvalue-sum bound is Ky Fan's variational principle,
using $|\Tr(PE)|\le\norm E_{\mathcal S_1}$ for finite-rank orthogonal
projections $P$ and trace-class $E$.
\end{proof}

Equation~\eqref{eq:layer} also proves directly that the positive operator
$K$ has infinite trace, since $\sum_p1/(p+1)$ diverges.
The trace-class perturbation statement alone is not used to infer
\eqref{eq:maineigen}; the mass scale here is below $1/r$, so a generic
trace-class perturbation does not supply the required relative control.

\section{Relation to the sparse finite matrices}\label{sec:finite}
This section recalls the finite matrices of~\cite{Kline} and the order of
limits in their relation to $K$. Proposition~\ref{prop:finite} is
\cite[Theorem~4.3]{Kline}; we include its proof so that the paper can be
read on its own. Corollary~\ref{cor:growing} records what an explicit
operator-norm estimate would add.

For $1\le i,j\le n$, let
\[
 (A_n)_{ij}=\one_{\{i=j\}}+
 \one_{\{i>1,\ i\in\sfq,\ j=i/P^+(i)\}},
 \qquad C_n=A_n^{-1}.
\]
The off-diagonal part is strictly lower triangular. Its powers follow
successive parent steps, so its finite geometric inverse gives
\begin{equation}\label{eq:Centries}
 (C_n)_{ij}=
 \begin{cases}
 \mu(i/j),&i,j\in\sfq,\ j\preceq i,\\
 1,&i=j\notin\sfq,\\
 0,&\text{otherwise}.
 \end{cases}
\end{equation}
In particular $C_ne_1=(\mu(1),\ldots,\mu(n))^{\mathsf T}$.

\begin{proposition}[Finite Gram convergence]\label{prop:finite}
Extend $H_n=C_n^*C_n/n$ by zero to $\ell^2(\N)$. Then
\[
 \norm{H_n-K}_{\HS}\longrightarrow0.
\]
Consequently, for every fixed $r\ge1$,
\[
 \frac{s_r(C_n)}{\sqrt n}\longrightarrow\sqrt{\lambda_r(K)},
\]
where $s_r(C_n)$ denotes the $r$th largest singular value.
\end{proposition}
\begin{proof}
For squarefree $j\preceq k\le n$, common descendant rows in
\eqref{eq:Centries} give
\[
 (C_n^*C_n)_{jk}=\mu(k/j)\,G_k(n),
 \quad
 G_k(n)=\#\{d\le n/k:\mu^2(d)=1,\ P^-(d)>P^+(k)\}.
\]
Here $P^-(1)=+\infty$. Incomparable squarefree coordinates give zero,
and nonsquarefree coordinates give only the diagonal entry $1$.

For fixed $y$, the density of squarefree integers coprime to
$P_y=\prod_{p\le y}p$ is
\begin{equation}\label{eq:roughdensity}
 \lim_{x\to\infty}\frac1x
   \sum_{\substack{d\le x\\(d,P_y)=1}}\mu^2(d)
 =\prod_{p\le y}(1-1/p)\prod_{p>y}(1-1/p^2)
 =c_0\prod_{p\le y}\frac p{p+1}.
\end{equation}
Indeed, substitute $\mu^2(d)=\sum_{\beta^2\mid d}\mu(\beta)$.
For each fixed $\beta$ coprime to $P_y$, the remaining coprimality
condition has density $\varphi(P_y)/P_y$; terms $\beta>B$, after division
by $x$, are bounded by $\sum_{\beta>B}\beta^{-2}$.
Taking $x\to\infty$ and then $B\to\infty$ gives the Euler product.
This proof includes $y=1$, with $P_y=1$.

It follows that $G_k(n)/n\to\delta_k$ for each fixed squarefree $k$.
Furthermore $G_k(n)/n\le1/k$ and $\delta_k\le1/k$. For maximal coordinate
$k$, there are at most $1+2\omega(k)$ ordered comparable pairs.
The convergent series
\[
 \sum_{k\ge1}\frac{1+2\omega(k)}{k^2}<\infty
\]
therefore dominates the squared squarefree matrix entries and their
differences, up to an absolute constant. Its convergence follows, for
example, from $\omega(k)\le\log k/\log2$ for $k\ge2$.
Dominated convergence gives the Hilbert--Schmidt limit on the squarefree
coordinates. The squared Hilbert--Schmidt norm on nonsquarefree
coordinates is at most $n/n^2=1/n$.
Finally, operator-norm convergence and the min--max principle imply
convergence of each fixed ordered eigenvalue of $H_n$. Taking square
roots proves the singular-value assertion.
\end{proof}

In particular, the rigorously justified iterated limit is
\[
 \lim_{r\to\infty}\sqrt{\frac{r\log r}{c_0}}\,
       \left(\lim_{n\to\infty}\frac{s_r(C_n)}{\sqrt n}\right)=1.
\]
A simultaneous limit requires a quantitative error estimate.

\begin{corollary}[A sufficient condition for growing indices]\label{cor:growing}
Suppose $\varepsilon_n\ge\norm{H_n-K}$ and $r_n\le n$ is an integer
sequence with $r_n\to\infty$ and
$\varepsilon_n r_n\log r_n\to0$. Then
\[
 s_{r_n}(C_n)\sim
       \sqrt{\frac{c_0 n}{r_n\log r_n}}.
\]
\end{corollary}
\begin{proof}
The min--max principle gives
$|\lambda_{r_n}(H_n)-\lambda_{r_n}(K)|\le\varepsilon_n$.
Divide by Theorem~\ref{thm:main}'s asymptotic for $\lambda_{r_n}(K)$,
then take square roots.
\end{proof}

\begin{remark}[Bordering and Mertens sums]
Let $\mathbf 1_n$ be the all-ones vector and set
\[
 B_n=A_n+e_1(\mathbf 1_n-e_1)^{\mathsf T}.
\]
Thus the first row is replaced by ones. The determinant lemma, together
with \eqref{eq:Centries}, gives the exact identity
\[
 \det B_n
 =1+(\mathbf 1_n-e_1)^{\mathsf T}C_ne_1
 =\sum_{j\le n}\mu(j)=M(n).
\]
The present counting theorem is unconditional and concerns the positive
limiting Gram spectrum. This determinant identity does not convert it into
a new cancellation estimate for $M(n)$. Such an application needs further
control of the signed vector and the bordering operation, or uniform
finite-size estimates strong enough for the intended arithmetic conclusion.
\end{remark}

\section{Related work and open questions}\label{sec:related}
\subsection{Related work}
\emph{Arithmetic matrices.} The fixed-index part of the connection in
Section~\ref{sec:finite} follows a pattern used by Hilberdink~\cite{Hil17}
for multiplicative Toeplitz matrices $(f(i/j)\one_{\{j\mid i\}})$: after
normalization, Gram matrices of the finite matrices converge to a positive
Hilbert--Schmidt operator, so each fixed singular value has a limit
described by an eigenvalue of that operator. Hilberdink and
Pushnitski~\cite{HP21,HP24} then determined eigenvalue asymptotics for
infinite arithmetic matrices built from least common multiples, such as
$(n^\sigma m^\sigma/[n,m]^\tau)$, obtaining power laws such as
$\lambda_n\sim\kappa n^{-\rho}$. Their proofs use a factorization of these
matrices as infinite tensor products over primes, which expresses the
eigenvalues as products of local eigenvalues, followed by a Tauberian
theorem of Wiener--Ikehara type. They also show that the constant $\kappa$
need not be the one suggested by the diagonal entries; see~(1.5) of
\cite{HP21} and the discussion following it. The kernel $K$ is of a
different kind: its entries depend
on the ancestor relation in the tree, not on a multiplicative function of
the indices, and we have not found a tensor-product factorization of it.
Its eigenvalue law also carries a logarithmic factor rather than being a
pure power. In contrast with the example of~\cite{HP21}, for $K$ the
leading constant is the one given by the count of diagonal masses.

\emph{Operators on trees.} The synthesis map $V$ of
\eqref{eq:gram-identification} is a discrete Hardy-type operator: for
finitely supported $x$, $(Vx)(\omega)$ is the sum of $x_d$ over the
vertices $d$ on the path determined by $\omega$. Evans, Harris and
Lang~\cite{EHL} proved, for Hardy operators on metric trees with finite
vertex degrees and suitable weights $u,v$, that the approximation numbers
satisfy $n\,a_n\to\alpha_p\int|uv|$, extending earlier Hilbert-space
results cited there. When this limit is finite,
$a_n=O(1/n)$. For $V$, Theorem~\ref{thm:main} and
\eqref{eq:gram-identification} give
$a_n(V)=(\lambda_n(K)/c_0)^{1/2}\sim(n\log n)^{-1/2}$, so $V$ is not covered
by that theorem with a finite integral. The tree of
Section~\ref{sec:kernel} has vertices of infinite degree, but this is a
feature of the presentation rather than of the operator. Revealing the
primes one at a time gives a binary filtration of $\Omega$ whose cylinders
$B$ are determined by the selected primes up to some prime $p$. Each $A_d$
is the cylinder of this kind with $p=P^+(d)$ in which the last revealed
prime is selected, and
\[
 T=\sum_B c_B\,\nu(B)^{-1}\rankone{\one_B},
 \qquad
 c_B=\begin{cases}\nu(B),&B=\Omega\text{ or the last revealed prime of }B
 \text{ is selected},\\0,&\text{otherwise}.\end{cases}
\]
Results on weighted dyadic or martingale Carleson embeddings, and on
singular values of weighted operators on trees of bounded degree, are
therefore relevant comparisons. We have not determined whether any of them
yields the exact constant in Theorem~\ref{thm:main}, only two-sided
bounds, or neither.

\emph{Ultrametric spaces.} For almost every pair $(\omega,\omega')$, the
kernel of $T$ is one plus the length of the longest common initial segment
of the increasing sequences of primes selected by $\omega$ and $\omega'$,
equivalently, the number of vertices whose cylinders contain both points.
This is an ultrametric quantity. Bendikov, Grigor'yan, Pittet and
Woess~\cite{BGPW} studied isotropic Laplacians on ultrametric spaces with
compact balls and a reference measure, which we again call $\nu$; their
eigenfunctions are the functions $\one_C/\nu(C)-\one_B/\nu(B)$, for a ball
$C$ with parent ball $B$~\cite[Theorem~3.8]{BGPW}. The operator $T$ is not diagonal in the
analogous basis. For $h=3\one_{A_2}-1$, which is this function for the
cylinder $A_2$ and its parent $\Omega$, one has $\inner{\one}{h}=0$ and
$(Th)(\omega)=2\sum\nu_d$ for $\omega\in A_2$, the sum taken over the
vertices $d\ne1$ on the path of $\omega$. This is not constant on $A_2$:
on $A_6$ it is at least $2\nu_2+2\nu_6=2/3+1/6$, while on $A_{10}$ the
vertices after $2$ have total mass at most $\nu_{10}/(1-1/7)=7/144<1/12$,
because every later child ratio is at most $1/7$ by \eqref{eq:arith-ratios};
so $Th<2/3+1/6$ there. Hierarchical Schr\"odinger operators, which add a potential
to a hierarchical Laplacian, are studied by Bendikov, Grigor'yan and
Molchanov~\cite{BGM}, including questions about discrete spectrum. Their
stated setting is homogeneous, and we have not carried out a comparison
with the inhomogeneous operator $T$.

\emph{Self-similar spectral problems.} The recursion behind
Proposition~\ref{prop:sandwich}, in which a subtree operator is a rank-one
operator plus rescaled copies of smaller subtree operators, resembles the
recursive counting identities used for spectral asymptotics on
self-similar sets, where renewal theory supplies the asymptotics; see
Kigami and Lapidus~\cite{KL}. Appendix~\ref{sec:stationary} uses the key
renewal theorem in exactly this way for stationary trees. For the
arithmetic tree the subtrees are not rescaled copies of a single operator,
and the asymptotics come instead from comparison with vertex counts.

\emph{The arithmetic count.} Theorem~\ref{thm:arithmetic} is an
application of classical estimates for smooth numbers~\cite{HT}, of the
same kind as those used for sums involving the largest prime factor; for
instance, Erd\H{o}s, Ivi\'c and Pomerance~\cite{EIP} evaluate
$\sum_{n\le x}1/P^+(n)$ through $\Psi(x,y)$ and $\rho$. We do not regard
the count as new in method.

We have not found Theorem~\ref{thm:main} in the sources discussed above.
The comparisons described there as not carried out could change that
assessment.

\subsection{Open questions}
\begin{question}\label{q:drop}
Under \eqref{eq:tree-assumptions}, suppose $D$ is regularly varying at
zero with index $-1$. Must $\ncount(T;t)\sim D(t)$ hold without either
vanishing child ratios or the profile hypotheses of
Theorem~\ref{thm:profile}?
\end{question}

The exact decomposition \eqref{eq:exact-shifts} reduces this question to
the locations of unit-integral, zero--one profiles at varying mass scales.
Their common integral does not, by itself, justify exchanging a limit
and the vertex sum. Proposition~\ref{prop:integrated} gives an integrated
comparison without profile hypotheses, but does not establish a pointwise
counting asymptotic. Proposition~\ref{prop:kary} is not a counterexample to
Question~\ref{q:drop}, because it lacks the stipulated regular variation.

Trees whose splitting rule depends on the position of a vertex's mass on
a logarithmic scale are natural candidates for a negative answer. For such
trees we have proved neither that the mass count is regularly varying nor
that $\ncount(T;t)/D(t)$ fails to converge, and we make no claim about
them here.

\begin{question}
Can the proof of Theorem~\ref{thm:arithmetic} and the sandwich in
Proposition~\ref{prop:sandwich} be made quantitative enough to give a
useful remainder in \eqref{eq:maincount}?
\end{question}
The proof as written takes $X\to\infty$ for fixed $a,b$, then sends
$a\downarrow0$ and $b\uparrow1$. The compact-$u$ input and this order
of limits yield a leading equivalent only. In particular, the present
argument does not identify a second term or a rate.

\begin{question}
Can one bound $\norm{C_n^*C_n/n-K}$ explicitly in a range that makes
Corollary~\ref{cor:growing} effective for useful $r_n$?
\end{question}
Hilbert--Schmidt convergence suffices for fixed indices, but supplies no
explicit scale in the proof given here. A quantitative treatment must
control both the rough-number densities and the matrix tail, then compare
that error with $1/(r_n\log r_n)$.

\begin{remark}[Arithmetic reach]
The main result concerns a positive kernel whose signs can be removed by
a diagonal unitary. Its proof uses classical smooth-number estimates and
does not produce new cancellation in M\"obius sums. An application to a
zero-free region would require an additional theorem that controls the
relevant signed arithmetic quantity; the asymptotic
$\lambda_r(K)\sim c_0/(r\log r)$ is not such a theorem by itself.
\end{remark}

\begin{remark}[Attribution]
The compactness argument, interlacing, the rank-one trace identity, the
renewal theorem, and the smooth-number inputs belong to established
operator theory, probability, and analytic number theory. We make no
priority claim for the abstract tree results in Section~\ref{sec:tree}
or in the appendices.
\end{remark}

\subsection*{Role of AI}
AI assisted the research development, literature search, proof exposition,
and mathematical checking under the author's direction. Process-separated
AI reviews are not independent expert peer review, and no independent human
expert has examined the proofs. The release repository records the
inspected sources, the proof reviews, and their dispositions.

\appendix
\section{Spectral shifts and a weaker transfer hypothesis}\label{sec:profiles}
This appendix gives an alternative sufficient condition for transfer.
It is not needed for Theorem~\ref{thm:main}. The tree assumptions are
\eqref{eq:tree-assumptions} throughout.

\subsection{Unit-integral rank-one shifts}
\begin{lemma}\label{lem:rankone-shift}
Let $B$ be positive and compact on a Hilbert space, and let $P=R_\phi$.
For $t>0$, put
\[
 \xi(t)=\ncount(B+P;t)-\ncount(B;t).
\]
Then $\xi(t)\in\{0,1\}$, it vanishes for $t\ge\norm{B+P}$, and
\begin{equation}\label{eq:shift-integral}
 \int_0^\infty\xi(t)\,dt=\norm\phi^2.
\end{equation}
No trace-class assumption on $B$ is required.
\end{lemma}
\begin{proof}
Positivity and rank-one interlacing give $0\le\xi\le1$. The support bound
is immediate. For $s>0$, both positive parts below have finite rank and
the layer-cake formula gives
\[
 g(s):=\Tr(B+P-sI)_+-\Tr(B-sI)_+
      =\int_s^\infty\xi(t)\,dt.
\]
If $F$ is the spectral projection of $B+P$ on $(s,\infty)$, the
variational principle for the trace of the positive part gives
\[
 g(s)\le\Tr(FP)\le\norm\phi^2.
\]
For the reverse bound, let $E=\one_{(s,\infty)}(B)$ and
$\psi=(I-E)\phi$. If $\psi=0$, the trial projection $E$ gives
$g(s)\ge\norm\phi^2$. Otherwise take the trial projection
$Q=E+R_{\psi/\norm\psi}$. Since $Q\phi=\phi$,
\[
 g(s)\ge\norm\phi^2+
    \inner{\psi/\norm\psi}{B(\psi/\norm\psi)}-s
 \ge\norm\phi^2-s.
\]
Let $s\downarrow0$ and use monotone convergence of the nonnegative
integral.
\end{proof}

For each vertex set
\[
 \widehat B_v=\widehat T_v-R_{\one},\qquad
 \widehat\xi_v(u)=
   \ncount(\widehat T_v;u)-\ncount(\widehat B_v;u).
\]
Then
\begin{equation}\label{eq:profile-facts}
 \widehat\xi_v\in\{0,1\},\qquad
 \widehat\xi_v(u)=0\ (u\ge M_\theta),\qquad
 \int_0^\infty\widehat\xi_v(u)\,du=1.
\end{equation}

\begin{proposition}[Exact counting decomposition]\label{prop:exact-shifts}
For every $t>0$,
\begin{equation}\label{eq:exact-shifts}
 \ncount(T;t)=\sum_v\widehat\xi_v(t/m_v),
 \qquad
 D(t)=\sum_v\one_{\{t/m_v<1\}}.
\end{equation}
Only finitely many terms in either sum can be nonzero.
\end{proposition}
\begin{proof}
The decomposition \eqref{eq:normalized-decomposition} gives
\[
 \ncount(T_v;t)=\widehat\xi_v(t/m_v)
                 +\sum_{c\text{ child of }v}\ncount(T_c;t).
\]
Unroll this identity from the root. At depth $h$, every subtree norm is
at most $M_\theta\theta^h$, so all residual counts vanish once this is
at most $t$. Nonzero shifts require $m_v>t/M_\theta$, a finite set.
The identity for $D$ is its definition.
\end{proof}

\begin{proposition}[An integrated comparison]\label{prop:integrated}
For every $t>0$,
\begin{equation}\label{eq:integrated-bound}
 0\le\int_t^\infty\{\ncount(T;s)-D(s)\}\,ds
 \le tD(t/M_\theta).
\end{equation}
\end{proposition}
\begin{proof}
For each vertex the function
$f_v(s)=\widehat\xi_v(s/m_v)$ takes values in $[0,1]$, is supported in
$[0,M_\theta m_v]$, and has integral $m_v$.
If $t<m_v$, then
\[
 \int_t^\infty(f_v(s)-\one_{\{s<m_v\}})\,ds
 =t-\int_0^t f_v(s)\,ds\in[0,t].
\]
If $t\ge m_v$, the same expression is $\int_t^\infty f_v(s)\,ds$,
which lies in $[0,m_v]\subseteq[0,t]$. It is zero when
$t\ge M_\theta m_v$. Sum over the finite set of possible contributors
using \eqref{eq:exact-shifts}.
\end{proof}

\subsection{Transfer through limiting profiles}
The following precise domination hypothesis is stronger than pointwise
convergence of the profiles. It states what is needed to exchange the
varying vertex sum with a limiting profile.

\begin{theorem}[Profile transfer at index $-1$]\label{thm:profile}
Suppose $D$ is regularly varying at zero with index $-1$.
Partition the vertices into an exceptional class $\mathcal E$ and
finitely many classes $\mathcal V_1,\ldots,\mathcal V_q$. Define
\[
 D_{\mathcal U}(t)=\#\{v\in\mathcal U:m_v>t\}.
\]
Assume
\[
 D_{\mathcal E}(t)=o(D(t)),\qquad
 \frac{D_{\mathcal V_k}(t)}{D(t)}\longrightarrow\beta_k>0,
 \qquad \sum_{k=1}^q\beta_k=1.
\]
For each $k$, suppose there is a bounded Riemann-integrable function
$h_k$ on $[0,M_\theta]$, extended by zero outside that interval.
For every $\varepsilon>0$, suppose there is a nonnegative bounded
Riemann-integrable function $F_{k,\varepsilon}$ on the same interval,
also extended by zero, such that
\[
 \int_0^{M_\theta}F_{k,\varepsilon}(u)\,du\longrightarrow0
 \quad(\varepsilon\downarrow0),
\]
and, outside a finite subset of $\mathcal V_k$ depending on $\varepsilon$,
\begin{equation}\label{eq:profile-domination}
 |\widehat\xi_v(u)-h_k(u)|\le F_{k,\varepsilon}(u)
 \qquad\text{for every }u>0.
\end{equation}
Then $\ncount(T;t)\sim D(t)$.
\end{theorem}
\begin{proof}
We first record a weighted-count consequence of regular variation.
If $\mathcal U=\mathcal V_k$ and $g$ is bounded and Riemann integrable
on $[0,M_\theta]$, extended by zero, then
\begin{equation}\label{eq:weighted-count}
 \frac1{D(t)}\sum_{v\in\mathcal U}g(t/m_v)
 \longrightarrow\beta_k\int_0^{M_\theta}g(u)\,du.
\end{equation}
To see this, the normalized counting measure of the points $t/m_v$
assigns mass
\[
 \frac{D_{\mathcal U}(t/b)}{D(t)}\longrightarrow\beta_k b
\]
to $(0,b)$, for each $b>0$. Its mass at a fixed positive endpoint tends
to zero: bracket that mass by the difference between the counts at
$b-\delta$ and $b+\delta$, take the limit, and let $\delta\downarrow0$.
Mass near zero is bounded by the mass of $(0,\delta)$, whose limit is
$\beta_k\delta$. Step-function upper and lower bounds, or Darboux sums
on a fixed finite partition, now prove \eqref{eq:weighted-count}.
This also accounts for possible atoms at $M_\theta$.

Each class $\mathcal V_k$ is infinite. Integrating
\eqref{eq:profile-domination} at any vertex outside its finite exception
and using \eqref{eq:profile-facts} gives
\[
 \left|1-\int h_k\right|\le\int F_{k,\varepsilon}.
\]
Thus $\int h_k=1$. For fixed $\varepsilon$, the finitely many exceptional
vertices in each class contribute $O(1)$ after subtracting the
corresponding $h_k$ terms. Apply \eqref{eq:weighted-count} to both
$h_k$ and $F_{k,\varepsilon}$ in \eqref{eq:profile-domination}. The
upper limiting normalized error is at most
$\beta_k\int F_{k,\varepsilon}$.

Finally the entire class $\mathcal E$ contributes at most
$D_{\mathcal E}(t/M_\theta)=o(D(t))$, by regular variation at the fixed
scale $1/M_\theta$. Sum over $k$ and let $\varepsilon\downarrow0$.
Equation~\eqref{eq:exact-shifts} then gives the limit $\sum_k\beta_k=1$.
\end{proof}

\begin{remark}[Recovery of vanishing child ratios]
When $a_v\le\eta$, \eqref{eq:normalized-decomposition} gives
\[
 \norm{\widehat B_v}\le M_\theta\eta,\qquad
 1\le\lambda_1(\widehat T_v)\le1+M_\theta\eta,\qquad
 \lambda_2(\widehat T_v)\le M_\theta\eta.
\]
Consequently $\widehat\xi_v(u)=1$ for $M_\theta\eta<u<1$, and it is
zero for $u\ge1+M_\theta\eta$. Thus the common profile
$h=\one_{[0,1)}$ works in Theorem~\ref{thm:profile}, with error envelope
\[
 F_\eta=\one_{[0,M_\theta\eta]}
          +\one_{[1,1+M_\theta\eta]}
\]
restricted to $[0,M_\theta]$, for sufficiently small $\eta$.
Its integral is at most $2M_\theta\eta$.
This recovers the index-$-1$ case of Theorem~\ref{thm:transfer}.
The argument uses only the upper bound
$\norm{\widehat T_v}-1\le\norm{\widehat B_v}\le M_\theta\eta$;
no lower bound on $\norm{\widehat T_v}-1$ is needed.
\end{remark}

\begin{remark}
At index $-\alpha$, the measure in \eqref{eq:weighted-count} would instead
be $\beta_k\alpha u^{\alpha-1}\,du$. The unit-integral identity
\eqref{eq:profile-facts} fixes its contribution without additional moments
precisely when $\alpha=1$. The profile theorem is therefore stated only
at that index.
\end{remark}

\section{Stationary trees and a lattice comparison}\label{sec:stationary}
\begin{proposition}[Finite stationary non-arithmetic trees]\label{prop:stationary}
Fix $k\ge2$ and $\rho_1,\ldots,\rho_k>0$ with $\sum_i\rho_i=1$.
On $\{1,\ldots,k\}^{\N}$ with the corresponding product probability,
take the cylinders indexed by all finite words, including the empty word.
Let $T$ be their cylinder operator. Put
\[
 a_i=-\log\rho_i,\qquad H=\sum_{i=1}^k\rho_i a_i.
\]
Assume there is no $h>0$ such that every $a_i$ belongs to $h\mathbb Z$.
Then, as $t\downarrow0$,
\[
 D(t)\sim\ncount(T;t)\sim\frac1{Ht}.
\]
\end{proposition}
\begin{proof}
The masses are $m_w=\prod_{j=1}^{|w|}\rho_{w_j}$, and
$\theta=\max_i\rho_i<1$. Vertices with mass above $t$ have bounded depth,
so $D(t)<\infty$. Proposition~\ref{prop:compact} applies.

Let $F$ be the probability measure assigning mass $\rho_i$ to $a_i$,
combining equal atoms if necessary, and let
$U=\sum_{\ell\ge0}F^{*\ell}$ be its renewal measure, including the
unit atom at zero. For $g(x)=e^{-x}\one_{\{x>0\}}$ and
$z(x)=e^{-x}D(e^{-x})$, direct summation over words gives
\[
 z(x)=\int_{[0,\infty)}g(x-y)\,dU(y).
\]
Indeed, a word contributes probability $m_w=e^{-\sum_j a_{w_j}}$
to the renewal measure, cancelling the factor $e^{\sum_j a_{w_j}}$
in $g(x-\sum_j a_{w_j})$.

The increments have finite positive mean $H$ and are non-arithmetic
by hypothesis. The function $g$ is directly Riemann integrable.
The key renewal theorem therefore gives
\[
 z(x)\longrightarrow\frac1H\int_0^\infty e^{-u}\,du=\frac1H.
\]
For a precise source, see Gallager~\cite{Gallager}, Theorem~4.7.3,
printed page~202. That theorem uses the renewal measure without the
initial atom; the additional term here is $g(x)\to0$.
Thus $D(t)\sim1/(Ht)$.

All normalized subtrees are unitarily equivalent. Their shifts
$\widehat\xi_v$ are consequently one common function $h$.
It is bounded by one and supported in $[0,M_\theta]$, and has integral
one by Lemma~\ref{lem:rankone-shift}. Away from zero its discontinuities
are among finitely many eigenvalues on each compact subinterval; it is
therefore Riemann integrable on $[0,M_\theta]$. Apply
Theorem~\ref{thm:profile} with one class and zero error envelopes to get
$\ncount(T;t)\sim D(t)$.
\end{proof}

For example, the weights $(3/10,7/10)$ satisfy the non-arithmetic
condition. A rational relation between $\log(10/3)$ and $\log(10/7)$
would, after exponentiation and comparison of prime exponents, give a
nontrivial equality between powers involving distinct primes.

\begin{proposition}[Homogeneous $k$-ary trees]\label{prop:kary}
In the same construction with $\rho_i=1/k$ for all $i$, the spectrum is
\[
 \frac{k}{k-1}\quad\text{with multiplicity }1,\qquad
 \frac{k^{-\ell}}{k-1}\quad\text{with multiplicity }(k-1)k^\ell
 \quad(\ell\ge0).
\]
In particular,
\begin{equation}\label{eq:kary-count}
 \ncount(T;t)=\one_{\{t<k/(k-1)\}}+(k-1)D((k-1)t),
 \qquad
 D(t)=\sum_{\ell\ge0}k^\ell\one_{\{k^{-\ell}>t\}}.
\end{equation}
The function $D$ is not regularly varying at zero with index $-1$.
For $k=3$,
\[
 \liminf_{t\downarrow0}\frac{\ncount(T;t)}{D(t)}=\frac23,\qquad
 \limsup_{t\downarrow0}\frac{\ncount(T;t)}{D(t)}=2.
\]
\end{proposition}
\begin{proof}
Let $E_\ell$ be conditional expectation onto the first $\ell$ coordinates,
with $E_0$ the projection onto constants. The sum of cylinder rank-one
operators at depth $\ell$ is $k^{-\ell}E_\ell$. Thus
$T=\sum_{\ell\ge0}k^{-\ell}E_\ell$.
Constants have eigenvalue $\sum_{\ell\ge0}k^{-\ell}=k/(k-1)$.
The martingale difference space
$\operatorname{ran}(E_{\ell+1}-E_\ell)$ has dimension $(k-1)k^\ell$,
and $T$ acts on it by
$\sum_{j\ge\ell+1}k^{-j}=k^{-\ell}/(k-1)$.
These orthogonal spaces and the constants exhaust $L^2$.
This proves the spectrum and \eqref{eq:kary-count}.

For any fixed $c\in(1,k)$ and $t_m=k^{-m}$, the strict counting
convention gives $D(ct_m)=D(t_m)$. This is inconsistent with the ratio
$c^{-1}$ required for regular variation of index $-1$.

For $k=3$, in each mass interval the argument $2t$ crosses exactly one
mass level. The ratios $D(2t)/D(t)$ have limiting values $1/3$ and $1$,
by summing the geometric series in \eqref{eq:kary-count}.
The single constant eigenvalue is negligible as $t\downarrow0$,
giving the two stated limits.
\end{proof}

\begin{remark}
For $k=2$, \eqref{eq:kary-count} instead gives
$\ncount(T;t)=D(t)+1$ for $0<t<2$.
Thus a nonvanishing child ratio can coexist with exact asymptotic transfer.
The $k=3$ example shows that $\theta<1$ alone does not force transfer,
but it does not answer the question with regular variation retained:
its $D$ fails that hypothesis.
\end{remark}

\begin{thebibliography}{99}
\bibitem{BGM}
A.~Bendikov, A.~Grigor'yan and S.~Molchanov,
\emph{On the spectrum of the hierarchical Schr\"odinger operator},
preprint, 2018.
\href{https://arxiv.org/abs/1811.05210}{arXiv:1811.05210}.

\bibitem{BGPW}
A.~Bendikov, A.~Grigor'yan, Ch.~Pittet and W.~Woess,
\emph{Isotropic Markov semigroups on ultra-metric spaces},
Russian Mathematical Surveys \textbf{69} (2014), no.~4, 589--680.
\href{https://doi.org/10.1070/RM2014v069n04ABEH004907}{doi:10.1070/RM2014v069n04ABEH004907}.
Theorem numbers refer to the version read,
\href{https://arxiv.org/abs/1304.6271}{arXiv:1304.6271}.

\bibitem{EIP}
P.~Erd\H{o}s, A.~Ivi\'c and C.~Pomerance,
\emph{On sums involving reciprocals of the largest prime factor of an
integer}, Glasnik Matemati\v{c}ki \textbf{21(41)} (1986), 283--300.

\bibitem{EHL}
W.~D.~Evans, D.~J.~Harris and J.~Lang,
\emph{The approximation numbers of Hardy-type operators on trees},
Proceedings of the London Mathematical Society \textbf{83} (2001), 390--418.
\href{https://doi.org/10.1112/plms/83.2.390}{doi:10.1112/plms/83.2.390}.
Version read: \href{https://arxiv.org/abs/math/0003215}{arXiv:math/0003215}.

\bibitem{Gallager}
R.~G.~Gallager,
\emph{Discrete Stochastic Processes}, Chapter~4, \emph{Renewal Processes},
MIT OpenCourseWare course copy, 2011.
\href{https://ocw.mit.edu/courses/6-262-discrete-stochastic-processes-spring-2011/931ffa0940899c27f34b71ad64fd2bb0_MIT6_262S11_chap04.pdf}{Chapter PDF}.
Theorem~4.7.3, printed p.~202.

\bibitem{Hil17}
T.~Hilberdink,
\emph{Singular values of multiplicative Toeplitz matrices},
Linear and Multilinear Algebra \textbf{65} (2017), 813--829.
\href{https://doi.org/10.1080/03081087.2016.1204978}{doi:10.1080/03081087.2016.1204978}.

\bibitem{HP21}
T.~Hilberdink and A.~Pushnitski,
\emph{Spectral asymptotics for a family of LCM matrices},
St.\ Petersburg Mathematical Journal \textbf{34} (2023), 463--481.
\href{https://doi.org/10.1090/spmj/1764}{doi:10.1090/spmj/1764}.
Equation numbers refer to the version read,
\href{https://arxiv.org/abs/2110.14323}{arXiv:2110.14323}.

\bibitem{HP24}
T.~Hilberdink and A.~Pushnitski,
\emph{Spectral asymptotics for a family of arithmetical matrices and
connection to Beurling primes}, preprint, 2024.
\href{https://arxiv.org/abs/2401.06892}{arXiv:2401.06892}.

\bibitem{HT}
A.~Hildebrand and G.~Tenenbaum,
\emph{Integers without large prime factors},
Journal de Th\'eorie des Nombres de Bordeaux \textbf{5} (1993), 411--484.
\href{https://numdam.org/item/JTNB_1993__5_2_411_0.pdf}{Original article}.
Relevant locations: printed p.~414 (smooth counts and the Dickman limit),
p.~422, equation (2.3), and p.~426, Lemma~2.6.

\bibitem{KL}
J.~Kigami and M.~L.~Lapidus,
\emph{Weyl's problem for the spectral distribution of Laplacians on
p.c.f.\ self-similar fractals},
Communications in Mathematical Physics \textbf{158} (1993), 93--125.

\bibitem{Kline19}
J.~Kline, \emph{A sparser matrix representation of the Mertens function},
Linear Algebra and its Applications \textbf{581} (2019), 354--366.
\href{https://doi.org/10.1016/j.laa.2019.07.021}{doi:10.1016/j.laa.2019.07.021}.

\bibitem{Kline}
J.~Kline,
\emph{The singular spectrum of a sparse Mertens matrix},
version 0.1.0, 2026.
\href{https://doi.org/10.5281/zenodo.23004980}{doi:10.5281/zenodo.23004980}.
\href{https://github.com/jeff-kline/sparse-mertens-singular-spectrum}{Source repository}.
\end{thebibliography}

\end{document}
